\begin{table}[H]
\centering
\caption{Espera por despeje y anticipación por condición de llegada.}
\label{tab:latencia_por_condicion}
\small
\begin{tabular}{llrrrrrr}
\toprule
\textbf{Enlace} & \textbf{Llegada} & \textbf{n} & \textbf{$T_{desp}$ med.} & \textbf{$T_{desp}$ p95} & \textbf{$A_v$ med.} & \textbf{$A_v$ p95} & \textbf{$A_v$ mín.} \\
\midrule
Laptop (3-nov) & Verde & 11 & -- & -- & 14,82 & 14,94 & 13,67 \\
Laptop (3-nov) & Ámbar/todo rojo & 10 & 1,81 & 4,00 & 12,76 & 13,69 & 9,71 \\
Laptop (3-nov) & Rojo general & 8 & -- & -- & 14,31 & 15,94 & 13,94 \\
\addlinespace
Nube, S1 (15-nov) & Verde & 14 & -- & -- & 13,92 & 14,89 & 12,81 \\
Nube, S1 (15-nov) & Rojo general & 22 & -- & -- & 14,70 & 15,62 & 13,72 \\
\addlinespace
Nube, S2 (20-nov, mañana) & Verde & 13 & -- & -- & 14,72 & 15,75 & 12,93 \\
Nube, S2 (20-nov, mañana) & Rojo general & 23 & -- & -- & 13,93 & 14,93 & 12,66 \\
\bottomrule
\end{tabular}
\par\vspace{3pt}{\footnotesize\raggedright\textit{Nota:} Espera por despeje $T_{desp}$ (RF-04) y anticipación $A_v$ (RF-11) por condición de llegada: la fase del controlador cuando la RSU decide sobre la solicitud. Ámbar y todo rojo se agrupan porque la RSU trata ambas igual (encola la solicitud). Escenarios de lazo y desvío combinados, en s. Elaboración propia.\par}
\end{table}
